\subsection{\texorpdfstring{第~\ref{sec:lagrangian-mechanics}~节}{Lagrange 力学}}\label{subsec:solutions-lagrangian-mechanics}

\begin{description}[font=\normalfont,leftmargin=2em,labelwidth=1.5em]
\RGsolution{ex:complete-proof-of-euler-lagrange-equations}{complete-proof-of-euler-lagrange-equations}
若 $c\in\mathcal C$ 是作用量的临界点, $(U,x^1,\ldots,x^n)$ 是局部坐标图, 且 $t\in(a,b)$ 满足 $c(t)\in U$, 则可取 $\varepsilon>0$, 使 $c([t-\varepsilon,t+\varepsilon])\subset U$. 考虑在 $[t-\varepsilon,t+\varepsilon]$ 之外与 $c$ 相同的变分, 可知 $c$ 在时间区间 $(t-\varepsilon,t+\varepsilon)$ 上必须满足此局部坐标图中的 Euler–Lagrange 方程. 因此作用量的任一临界点, 对所有满足 $c(t)\in U$ 的 $t\in(a,b)$, 都必须满足局部坐标图 $(U,x^1,\ldots,x^n)$ 中的 Euler–Lagrange 方程.

反之, 设 $c\in\mathcal C$ 在任意局部坐标图中都满足 Euler–Lagrange 方程. 在 $M$ 上引入辅助的 Riemann 度量, 并考虑以 $c([a,b])$ 中各点为中心的正规球. 因为 $c([a,b])$ 紧, 可选取这些球使其半径的下确界为正 (考虑 $c([a,b])$ 由全正规邻域组成的开覆盖). 由于 $c$ 的长度必为有限, 可选取其中有限多个球 $B_1,\ldots,B_N$ (它们是正规坐标的坐标邻域), 以及满足 $a<t_1<\cdots<t_{N-1}<b$ 的点 $t_1,\ldots,t_{N-1}$, 使 $c(a)\in B_1$, $c(b)\in B_N$, 且
\[
  c(t_i)\in B_i\cap B_{i+1}\quad(i=1,\ldots,N-1).
\]
对由 $\gamma:(-\varepsilon,\varepsilon)\times[a,b]\to M$ 给出的任一变分 $\gamma$, 重复 Theorem~\ref{thm:euler-lagrange-equations} 的证明中的计算, 有
\[
  \left.\frac{\mathrm{d}}{\mathrm{d}s}\right|_{s=0}
     \int_{t_{i-1}}^{t_i}L(\dot\gamma(s))\,\mathrm{d}t
    =(\mathbb FL)_{\dot c(t_i)}\left(\frac{\partial\gamma}{\partial s}(0,t_i)\right)
     -(\mathbb FL)_{\dot c(t_{i-1})}\left(\frac{\partial\gamma}{\partial s}(0,t_{i-1})\right)
\]
对 $i=2,\ldots,N-1$ 成立, 其中使用了 $c$ 在 $B_i$ 上满足 Euler–Lagrange 方程. 类似地,
\[
  \left.\frac{\mathrm{d}}{\mathrm{d}s}\right|_{s=0}
     \int_a^{t_1}L(\dot\gamma(s))\,\mathrm{d}t
   =(\mathbb FL)_{\dot c(t_1)}\left(\frac{\partial\gamma}{\partial s}(0,t_1)\right)
\]
且
\[
  \left.\frac{\mathrm{d}}{\mathrm{d}s}\right|_{s=0}
     \int_{t_{N-1}}^b L(\dot\gamma(s))\,\mathrm{d}t
   =-(\mathbb FL)_{\dot c(t_{N-1})}\left(\frac{\partial\gamma}{\partial s}(0,t_{N-1})\right).
\]
将这些等式相加, 最终得到
\[
  \left.\frac{\mathrm{d}}{\mathrm{d}s}\right|_{s=0}\int_a^b L(\dot\gamma(s))\,\mathrm{d}t=0.
\]

\RGsolution{ex:restricted-three-body-problem}{restricted-three-body-problem}
\begin{enumerate}[label=(\alph*),ref=\alph*,leftmargin=*]
\item 若把 $\mathbb{R}^2$ 与 $\mathbb{R}^3$ 中的 $z=0$ 平面等同, 则旋转参考系的角速度为 $\Omega=e_z$. 由 Exercise~\ref{ex:rigid-body}(\ref{ex:inertial-forces-in-rotating-frame}), 第三个粒子的运动方程为
\begin{align*}
  m(\ddot x,\ddot y,\ddot z)
  &=(F_x,F_y,F_z)-m\Omega\times(\Omega\times(x,y,z))\\
  &\quad-2m\Omega\times(\dot x,\dot y,\dot z)-m\dot\Omega\times(x,y,z)\\
  &=(F_x,F_y,F_z)+m(x,y,0)-2m(-\dot y,\dot x,0).
\end{align*}
这个方程的前两个分量正是
\[
  \begin{cases}
   \ddot x=\dfrac{F_x}{m}+x+2\dot y,\\[3pt]
   \ddot y=\dfrac{F_y}{m}+y-2\dot x.
  \end{cases}
\]
(第三个分量是 $m\ddot z=F_z$, 因此对在轨道平面内运动的粒子, 要求 $F_z=0$.)
\item 首先注意
\[
  r_1=((x-1+\mu)^2+y^2)^{\frac12}
  \quad\text{且}\quad r_2=((x+\mu)^2+y^2)^{\frac12}.
\]
于是,
\begin{align*}
  \frac{\partial L}{\partial x}&=v^y+x-\frac{\mu}{r_1^3}(x-1+\mu)-\frac{1-\mu}{r_2^3}(x+\mu),\\
  \frac{\partial L}{\partial y}&=-v^x+y-\frac{\mu}{r_1^3}y-\frac{1-\mu}{r_2^3}y.
\end{align*}
此外,
\[
  \frac{\partial L}{\partial v^x}=v^x-y
   \quad\text{且}\quad\frac{\partial L}{\partial v^y}=v^y+x,
\]
所以
\begin{align*}
  \frac{\mathrm{d}}{\mathrm{d}t}
    \left(\frac{\partial L}{\partial v^x}(x,y,\dot x,\dot y)\right)&=\ddot x-\dot y,\\
  \frac{\mathrm{d}}{\mathrm{d}t}
    \left(\frac{\partial L}{\partial v^y}(x,y,\dot x,\dot y)\right)&=\ddot y+\dot x.
\end{align*}
因此 Euler–Lagrange 方程为
\[
  \begin{cases}
   \ddot x-2\dot y=x-\dfrac{\mu}{r_1^3}(x-1+\mu)-\dfrac{1-\mu}{r_2^3}(x+\mu),\\[3pt]
   \ddot y+2\dot x=y-\dfrac{\mu}{r_1^3}y-\dfrac{1-\mu}{r_2^3}y,
  \end{cases}
\]
这恰好是旋转参考系中的运动方程.
\item Hamilton 函数为
\begin{align*}
  H&=v^x\frac{\partial L}{\partial v^x}+v^y\frac{\partial L}{\partial v^y}-L\\
   &=\frac12((v^x)^2+(v^y)^2)-\frac12(x^2+y^2)-\frac{\mu}{r_1}-\frac{1-\mu}{r_2}.
\end{align*}
\item 现在求平衡点, 即运动方程的常值解. 由于这时有
\[
  \dot x=\ddot x=\dot y=\ddot y=0,
\]
故得到
\[
  \begin{cases}
   x-\dfrac{\mu}{r_1^3}(x-1+\mu)-\dfrac{1-\mu}{r_2^3}(x+\mu)=0,\\[3pt]
   y\left(1-\dfrac{\mu}{r_1^3}-\dfrac{1-\mu}{r_2^3}\right)=0.
  \end{cases}
\]
若 $y\ne0$, 则由第二个方程得
\[
  1-\frac{\mu}{r_1^3}-\frac{1-\mu}{r_2^3}=0.
\]
将第一个方程写成
\[
  x\left(1-\frac{\mu}{r_1^3}-\frac{1-\mu}{r_2^3}\right)
   +\mu(1-\mu)\left(\frac1{r_1^3}-\frac1{r_2^3}\right)=0,
\]
再利用第二个方程, 可知在平衡点处 $r_1=r_2=1$, 这也满足第一个方程. 因此, 有一个平衡点 $(\frac12-\mu,\frac{\sqrt3}{2})$ 满足 $y>0$, 另一个平衡点 $(\frac12-\mu,-\frac{\sqrt3}{2})$ 满足 $y<0$. 注意, 这两点各自到两个有质量的粒子的距离相等.

若 $y=0$, 则平衡点由函数
\[
  U(x)=\frac12x^2+\frac{\mu}{|x+1-\mu|}+\frac{1-\mu}{|x+\mu|}
\]
的临界点给出. 由于
\[
  U''(x)=1+\frac{2\mu}{|x+1-\mu|^3}+\frac{2(1-\mu)}{|x+\mu|^3}>0,
\]
且
\[
  \lim_{x\to\pm\infty}U(x)=\lim_{x\to\mu}U(x)=\lim_{x\to\mu-1}U(x)=+\infty,
\]
可知 $U$ 恰有三个临界点, 它们是局部极小点, 分别位于区间 $(-\infty,1-\mu)$, $(1-\mu,\mu)$ 和 $(\mu,+\infty)$ 中.
\item 为将系统线性化, 作代换
\[
  \begin{cases}
    x=\frac12-\mu+\xi,\\[2pt]
    y=\pm\frac{\sqrt3}{2}+\eta,
  \end{cases}
\]
并注意, 在平衡点处,
\begin{align*}
  \frac{\partial r_1}{\partial x}&=\frac{x-1+\mu}{r_1}=-\frac12;&
  \frac{\partial r_1}{\partial y}&=\frac y{r_1}=\pm\frac{\sqrt3}{2};\\
  \frac{\partial r_2}{\partial x}&=\frac{x+\mu}{r_2}=\frac12;&
  \frac{\partial r_2}{\partial y}&=\frac y{r_2}=\pm\frac{\sqrt3}{2},
\end{align*}
所以
\[
  \frac1{r_1^3}=1+\frac32\xi\mp\frac{3\sqrt3}{2}\eta+\cdots
\]
且
\[
  \frac1{r_2^3}=1-\frac32\xi\mp\frac{3\sqrt3}{2}\eta+\cdots.
\]
代入运动方程得线性化系统
\[
  \left\{\begin{aligned}
  \ddot\xi-2\dot\eta
   &=-\mu\left(\frac32\xi\mp\frac{3\sqrt3}{2}\eta\right)\left(-\frac12\right)\\
   &\quad-(1-\mu)\left(-\frac32\xi\mp\frac{3\sqrt3}{2}\eta\right)\left(\frac12\right),\\
  \ddot\eta+2\dot\xi
   &=-\mu\left(\frac32\xi\mp\frac{3\sqrt3}{2}\eta\right)\left(\pm\frac{\sqrt3}{2}\right)\\
   &\quad-(1-\mu)\left(-\frac32\xi\mp\frac{3\sqrt3}{2}\eta\right)\left(\pm\frac{\sqrt3}{2}\right),
  \end{aligned}\right.
\]
或等价地,
\[
  \begin{cases}
    \ddot\xi-2\dot\eta=\dfrac34\xi\pm\dfrac{3\sqrt3}{4}(1-2\mu)\eta,\\[4pt]
    \ddot\eta+2\dot\xi=\pm\dfrac{3\sqrt3}{4}(1-2\mu)\xi+\dfrac94\eta.
  \end{cases}
\]
相应的一阶线性系统的矩阵为
\[
  A=\begin{bmatrix}
   0&0&1&0\\
   0&0&0&1\\
   \frac34&\pm\frac{3\sqrt3}{4}(1-2\mu)&0&2\\
   \pm\frac{3\sqrt3}{4}(1-2\mu)&\frac94&-2&0
  \end{bmatrix},
\]
其特征多项式为
\[
  \det(A-\lambda I)=\lambda^4+\lambda^2+\frac{27}{4}\mu(1-\mu).
\]
该多项式的根满足
\[
  \lambda^2=\frac{-1\pm\sqrt{1-27\mu(1-\mu)}}{2},
\]
因此除非根都是纯虚数, 否则至少有一个根的实部为正. 所以平衡点不稳定当且仅当
\[
  1-27\mu(1-\mu)<0
   \iff\frac{1-\frac{\sqrt{69}}9}{2}<\mu<\frac{1+\frac{\sqrt{69}}9}{2}.
\]
\end{enumerate}

\RGsolution{ex:noether-for-free-rigid-body}{noether-for-free-rigid-body}
\begin{enumerate}[label=(\alph*),ref=\alph*,leftmargin=*]
\item 给定 $S\in SO(3)$, 有
\[
  \left.\frac{\mathrm{d}}{\mathrm{d}t}\right|_{t=0}\exp(tB)S
  =BS=(\mathrm{d}R_S)_IB=(X^B)_S,
\]
其中 $X^B$ 是由 $B$ 确定的右不变向量场.
\item 自由刚体的 Lagrange 函数是
\[
  L(V)=\frac12\langle\!\langle V,V\rangle\!\rangle,
\]
且显然在 $SO(3)$ 下不变 (因为 $\langle\!\langle\cdot,\cdot\rangle\!\rangle$ 不变). 由于
\[
  (\mathbb FL)_V(W)=\langle\!\langle V,W\rangle\!\rangle,
\]
Noether 定理保证对任一 $B\in\mathfrak{so}(3)$, 量
\[
  J^B=(\mathbb FL)_{\dot S}(X^B)
      =\langle\!\langle\dot S,BS\rangle\!\rangle=\langle\!\langle SA,BS\rangle\!\rangle
\]
沿系统的运动守恒, 其中照常记 $\dot S=SA$. 令 $\Omega=\Omega(A)$, $\Sigma=\Omega(B)$, 有
\begin{align*}
  J^B&=\int_{\mathbb{R}^3}\langle SA\xi,BS\xi\rangle\,\mathrm{d}m
      =\int_{\mathbb{R}^3}\langle S(\Omega\times\xi),\Sigma\times(S\xi)\rangle\,\mathrm{d}m\\
     &=\int_{\mathbb{R}^3}\langle\Omega\times\xi,(S^t\Sigma)\times\xi\rangle\,\mathrm{d}m
      =\int_{\mathbb{R}^3}\langle(S^t\Sigma),\xi\times(\Omega\times\xi)\rangle\,\mathrm{d}m\\
     &=\left\langle S^t\Sigma,\int_{\mathbb{R}^3}[\xi\times(\Omega\times\xi)]\,\mathrm{d}m\right\rangle
      =\langle S^t\Sigma,P\rangle=\langle\Sigma,SP\rangle.
\end{align*}
因为 $B$, 从而 $\Sigma$ 是任意的, 所以向量 $p=SP$ 守恒.
\end{enumerate}
\end{description}
